\section{Модуль персонализированной калибровки}

\subsubsection{Задача калибровки}

Программный комплекс отслеживает положение тела пациента в кадре
веб-камеры с помощью MediaPipe Pose~\cite{mediapipe_pose}: по каждому
кадру определяются координаты опорных точек (плечи, запястья, нос,
глаза и др.) в пикселях изображения. На этих данных строится логика
упражнений:
\begin{itemize}
    \item для упражнений с захватом объектов~--- проверка попадания
    кисти в область цели (\texttt{catch\_ok});
    \item для упражнений с движением головы~--- проверка выполнения
    поворота или наклона (\texttt{head\_ok});
    \item для всех упражнений~--- размещение целей в зоне досягаемости
    рук (\texttt{spawn\_object}) и контроль расстояния до камеры
    (\texttt{check\_runtime}).
\end{itemize}

Координаты landmarks зависят не только от анатомии пациента, но и от
расстояния до камеры, положения в кадре и разрешения видеопотока.
Если использовать для всех пользователей одни и те же координаты
(фиксированная зона появления целей, один радиус захвата, одни пороги
поворота головы), упражнения оказываются:
\begin{itemize}
    \item слишком лёгкими для одних пациентов и недостижимыми для других;
    \item нестабильными при смене дистанции до экрана между сессиями;
    \item некорректными при упражнениях на шею, где важен
    индивидуальный объём движений, а не абсолютные пиксели.
\end{itemize}

\textbf{Задача калибровки}~--- один раз (или при смене условий) измерить
персональные параметры пациента и сохранить их в профиле
$\mathcal{P}_u$. В дальнейшем все проверки выполнения и генерация
целей выполняются относительно $\mathcal{P}_u$, а не абсолютных
координат экрана.

Требования к модулю калибровки:
\begin{enumerate}
    \item не использовать маркеры и дополнительное оборудование~---
    только веб-камера;
    \item учитывать досягаемость \emph{обеих} рук и диапазон движений
    головы (упражнения 1--9 используют руки, шею или оба канала);
    \item зафиксировать эталонное расстояние до камеры, чтобы сессии
    были сопоставимы;
    \item быть понятным пациенту без участия врача;
    \item сохранять профиль между запусками программы.
\end{enumerate}

\subsubsection{Анализ существующих подходов}

\emph{Фиксированные зоны без калибровки.}
Простейший вариант~--- задавать координаты целей как доли ширины и
высоты кадра (например, 20\%--80\% по $x$ и $y$). Математически
$\mathbf{z} = (f_x \cdot w,\; f_y \cdot h)$, где $f_x, f_y$~---
константы. Подход не требует подготовительного этапа, но не
учитывает рост, длину рук и дистанцию до камеры; в реабилитации
с ним часто получают систематическую ошибку оценки
выполнения~\cite{hou2025}.

\emph{Калибровка по одной базовой позе (T-pose).}
Распространённый в CV-приложениях для реабилитации протокол: пациент
поднимает руки в стороны, система фиксирует размах рук и центр
торса~\cite{smartreflex}. Формально сохраняются ширина плеч $sw$ и
координаты кистей; зона упражнения масштабируется от $sw$. Этого
достаточно для контроля дистанции и грубой зоны рук, но не задаёт
досягаемость отдельно вверх, вниз, влево и вправо и не калибрует
голову.

\emph{Калибровка по одному повторению с максимальной амплитудой.}
В системе верхней конечности на MediaPipe пациент выполняет одно
повторение упражнения с максимально возможной амплитудой; по нему
записываются предельные углы суставов и на их основе задаются
допустимые диапазоны последующих повторений~\cite{hou2025}. Подход
персонализирует нагрузку, но привязан к конкретному типу движения
и плохо переносится на набор из девяти разнородных упражнений
(захват цели, движения шеи, разная геометрия).

\emph{Нормализация на анатомический масштаб.}
Чтобы уменьшить зависимость от дистанции до камеры, координаты
нормируют на локальный масштаб тела~--- чаще всего на ширину плеч
$sw$ или длину торса~\cite{mediapipe_pose,hou2025}. Например,
смещение головы задаётся как $\Delta x = (x_{\mathrm{nose}} -
x_{\mathrm{torso}}) / sw$, а не в пикселях. Это стандартный приём
в markerless pose estimation, но сам по себе не заменяет измерение
индивидуальных пределов движения~--- только делает метрики
сопоставимыми при небольших изменениях дистанции.

\emph{Экспертная настройка порогов.}
В системах контроля движений (например, функциональные тесты на
MediaPipe) пороги срабатывания подбираются разработчиком по серии
пробных записей и затем фиксируются в коде~\cite{soma2023}. Метод
точен при известной установке камеры, но не переносится на домашние
условия без персонализации.

\emph{Маркерные системы и depth-камеры.}
Оптические системы с маркерами или RGB-D камеры дают метрически точную
калибровку рабочего объёма, но противоречат требованию домашнего
тренажёра «камера ноутбука + без датчиков», принятому в
разрабатываемом комплексе и в аналогах вроде VestAid~\cite{hovareshti2021}.

Сопоставление по требованиям задачи:
\begin{itemize}
    \item \textbf{T-pose}~--- минимальная калибровка, нет данных по
    голове и по направлениям досягаемости;
    \item \textbf{одно max-ROM повторение}~\cite{hou2025}~---
    персонализация есть, но узко под одно упражнение;
    \item \textbf{нормализация на $sw$}~\cite{mediapipe_pose}~---
    необходима, но недостаточна без измерения пределов;
    \item \textbf{фиксированные зоны}~--- не подходят для
    реабилитационного тренажёра с объективным контролем.
\end{itemize}

Для комплекса с девятью типами упражнений нужен протокол, который
объединяет фиксацию дистанции (как в T-pose), измерение досягаемости
рук по направлениям и калибровку головы с нормализацией на $sw$.

\subsubsection{Выбранный подход и реализация}

Реализован модуль \texttt{calibration/} с классом \texttt{Calibrator}
и структурой данных \texttt{Profile}. Калибровка выполняется один раз
перед первым упражнением в сессии (или до завершения протокола) и
сохраняется в файл \texttt{calibration\_profiles/user\_\{id\}.json}.

\subsubsubsection{Протокол калибровки}

Калибровка организована как конечный автомат из трёх этапов:
\[
A \;\rightarrow\; B \;\rightarrow\; C \;\rightarrow\; \mathcal{P}_u.
\]

\textbf{Этап~A~--- положение и дистанция.}
Пациент разводит руки в стороны на уровне плеч (аналог T-pose).
Система проверяет:
\begin{itemize}
    \item обе кисти и плечи видны в кадре;
    \item кисти не выходят за поля кадра;
    \item ширина плеч $sw$ в пикселях попадает в допустимый диапазон
    (слишком близко / далеко~--- подсказка «подойдите» / «отойдите»);
    \item руки разведены достаточно широко.
\end{itemize}
При удержании позы 3~с фиксируются: центр торса $(x_{\mathrm{torso}},
y_{\mathrm{torso}})$, эталонная $sw$, начальный размах
$\mathit{span\_left}$, $\mathit{span\_right}$ относительно торса.

\textbf{Этап~B~--- досягаемость рук.}
Восемь шагов: для правой и левой руки по четыре направления
(вверх, вниз, влево, вправо). На экране показывается целевой круг;
пациент тянется к нему. Для каждого шага записывается максимальное
расстояние от торса до кисти в заданном направлении (в пикселях):
\[
R_{\mathrm{hand}}^{d} = \max_t \bigl\| \mathbf{p}_{\mathrm{wrist}}(t)
- \mathbf{p}_{\mathrm{torso}} \bigr\|_{\mathrm{axis}=d},
\quad d \in \{\uparrow, \downarrow, \leftarrow, \rightarrow\}.
\]

Если пациент не дотягивается до \emph{идеальной} цели (заданной от
начального размаха), применяется \textbf{адаптивное смещение цели}:
координаты круга плавно подтягиваются к фактической позиции кисти,
но не ближе торса и не дальше идеала. Так измеряется реальная
досягаемость, а не недостижимая «эталонная» точка.

\textbf{Этап~C~--- диапазон движений головы.}
Шесть положений: нейтральное, поворот влево и вправо, наклон
вверх и вниз, наклон к плечу. Для каждого измеряются метрики головы,
нормированные на $sw$:
\[
\delta_x = \frac{x_{\mathrm{nose}} - x_{\mathrm{torso}}}{sw}, \quad
\delta_y = \frac{y_{\mathrm{nose}} - y_{\mathrm{torso}}}{sw}, \quad
\varphi_{\mathrm{eye}} = \angle(\mathrm{LE}, \mathrm{RE}),
\]
где LE, RE~--- углы глаз. Порог «достаточного» поворота задаётся
относительно нейтрального положения (например,
$|\delta_x - \delta_x^{0}| \geq \theta_{\mathrm{turn}}$).
Каждая поза удерживается 2--2,5~с; усреднённые значения входят
в профиль.

\subsubsubsection{Структура профиля $\mathcal{P}_u$}

Профиль \texttt{Profile} содержит:
\begin{itemize}
    \item идентификатор пользователя, размер кадра;
    \item $sw$, $(x_{\mathrm{torso}}, y_{\mathrm{torso}})$~---
    эталон положения;
    \item $R_{\mathrm{right}}^{\uparrow,\downarrow,\leftarrow,\rightarrow}$,
    $R_{\mathrm{left}}^{\uparrow,\downarrow,\leftarrow,\rightarrow}$~---
    досягаемость рук;
    \item $\delta_x, \delta_y, \varphi_{\mathrm{eye}}$ для нейтрали
    и крайних положений головы;
    \item $r_{\mathrm{catch}}$~--- радиус захвата цели
    ($r_{\mathrm{catch}} \approx 0{,}12 \cdot sw$).
\end{itemize}

Профиль сериализуется в JSON; при следующем запуске загружается
функцией \texttt{load\_profile(user\_id)}.

\subsubsubsection{Использование профиля в упражнениях}

После калибровки все потоки упражнений (\texttt{detect\_thread.py},
\texttt{ex\_2.py}--\texttt{ex\_9.py}) работают с $\mathcal{P}_u$:

\begin{enumerate}
    \item \textbf{Контроль дистанции.}
    На каждом кадре сравнивается текущая $sw$ с эталонной:
    если $|sw - sw_0| / sw_0 > 0{,}4$, упражнение ставится на паузу
    с подсказкой «подойдите» или «отойдите». Это удерживает
    сопоставимость сессий без повторной калибровки.

    \item \textbf{Компенсация смещения торса.}
    Если пациент слегка сдвинулся в кадре, смещение
    $(\Delta x_{\mathrm{torso}}, \Delta y_{\mathrm{torso}})$
    учитывается при размещении целей (\texttt{torso\_shift}).

    \item \textbf{Генерация целей.}
    Координаты новой цели выбираются случайно внутри
    персонального прямоугольника досягаемости (\texttt{spawn\_object}),
    уменьшенного на 12\% от края (запас от границы движения).

    \item \textbf{Проверка захвата.}
    Захват засчитывается, если расстояние от кисти до центра цели
    меньше $r_{\mathrm{catch}}$ (\texttt{catch\_ok}).

    \item \textbf{Проверка движений головы.}
    В упражнениях на шею (\texttt{head\_ok}) текущие
    $(\delta_x, \delta_y, \varphi_{\mathrm{eye}})$ сравниваются
    с профилем: для режима «нейтрально»~--- близость к нейтрали,
    для «влево / вправо / …»~--- достижение доли пути между
    нейтралью и калиброванным крайним положением.
\end{enumerate}

Учёт зеркального отображения кадра (\texttt{cv2.flip}): после отражения
правая рука пользователя соответствует landmarks MediaPipe
\texttt{LEFT\_*}*, левая~--- \texttt{RIGHT\_*}*; модуль
\texttt{user\_hand\_px} выполняет это сопоставление явно.

\subsubsubsection{Связь с анализом подходов}

Выбранный протокол сочетает элементы рассмотренных решений:
\begin{itemize}
    \item этап~A~--- T-pose и фиксация $sw$, как в
    протоколах~\cite{smartreflex,hovareshti2021};
    \item этап~B~--- измерение индивидуальной досягаемости по
    направлениям, развивая идею max-ROM калибровки~\cite{hou2025}
    для набора из восьми осей, а не одного повторения;
    \item этап~C и функция \texttt{head\_ok}~--- нормализация на
    $sw$~\cite{mediapipe_pose} плюс запись \emph{личных} пределов
    шеи, чего нет в T-pose-only системах;
    \item адаптивное смещение цели на этапе~B~--- учёт ограниченной
    подвижности в реабилитации без отказа от калибровки;
    \item \texttt{check\_runtime}~--- поддержка стабильности между
    сессиями без маркерного оборудования.
\end{itemize}

Модуль калибровки связан с модулем рекомендаций параметров
разделением задач: калибровка задаёт \emph{геометрию} контроля
(где может быть цель и что считать выполненным), рекомендации~---
\emph{нагрузку} (сколько целей, с какой скоростью и в какие
интервалы). Оба модуля используют профиль пациента, но не пересекаются
по функциям.

Интеграция в интерфейс: калибровка запускается автоматически в
видеопотоке упражнения до начала игровой фазы; на кадре отображаются
подсказки, прогресс-бар и целевые маркеры. После успешного завершения
упражнение переходит к основному сценарию без отдельного экрана
настроек.
